% /chapters/1-intro/introduction.tex

\chapter{Introduction}\label{ch:introduction}

Consider the following version of the perfect set theorem.
\begin{theorem}[Cantor]\label{thm:cantor-perfect}
    Let $C$ be a closed subset of the Cantor space $2^\omega$.
    Then $C$ is uncountable if and only if there is a continuous injection from $2^\omega$ into $C$.
\end{theorem}
The statement says that the question of whether a closed subset of $2^\omega$ is uncountable has a structurally simple answer: one either embeds the Cantor space itself, or the set must be countable.
The Cantor space is the only obstruction to countability.

In the study of function spaces, a phenomenon of similar spirit manifests when comparing the sequential and topological closures of a set of functions.

On the Cantor space $X = 2^\omega$, the $n$-th coordinate map $p_n\colon 2^\omega\to2$ is continuous for each $n<\omega$; however, $(p_n)_{n\in\mathbb{N}}$ has \emph{no} convergent subsequences in $\mathbb{R}^X$ (\cite[Chapter~1.1]{Todorcevic_1997_TopicsTop}).
In a sense, this example exhibits the worst failure of sequential convergence possible: the closure of $\{p_n\}$ in $\{0, 1\}^X$ (or in $\mathbb{R}^X$ for that matter) is homeomorphic to the Stone-Čech compactification $\beta\mathbb{N}$ of the (discrete) space of natural numbers, which has cardinality $2^{2^{\omega}}$.

The following theorem, proved by Haskell Rosenthal in 1974, is fundamental in functional analysis and captures a sharp division in the behavior of sequences in a Banach space.
We denote by $\Cp(X)$ the space of continuous real-valued functions on $X$ with the topology of pointwise convergence.
Recall also that $\beta\mathbb{N}$ is the Stone-\v{C}ech compactification of the discrete space of natural numbers.

\begin{restatable}[Rosenthal's Dichotomy, \cite{Rosenthal:1974}]{theorem}{rosenthaldichotomy}\label{thm:Rosenthal}
    If $X$ is Polish and $(f_n)_{n\in\mathbb{N}}\subseteq \Cp(X)$ is pointwise bounded, then $(f_n)$ has a convergent subsequence, or a subsequence whose closure in $\mathbb{R}^X$ is homeomorphic to~$\beta\mathbb{N}$.
\end{restatable}

This result is a precursor to the classification theorems about Rosenthal compacta that we use in Chapter~\ref{ch:deep_computations} to establish new definability classes for various spaces of computations.

Let $G$ be a graph with vertex set $V(G)$ and edge set $E(G)$.
A $k$-colouring is a mapping from $V(G)$ to a set of $k$ colours such that adjacent vertices are given different colours.
We say $G$ is $k$-colourable if it admits a $k$-colouring.
The least number $k$ of colours for which a $k$-colouring of $G$ exists is denoted by $\chi(G)$.
Throughout this thesis, we will often endow the sets of vertices and edges with a topology and impose additional restrictions on the colouring, such as it being a Borel mapping.
In that case, the least number of colours $k$ for which there is a Borel $k$-colouring is denoted by $\chi_B(G)$.

Graph colourings are special instances of \emph{graph homomorphisms} between two graphs $f\colon G\to H$, that is, a mapping between the vertex sets that preserves the adjacencies: whenever $x,y$ form an edge in $G$, $f(x),f(y)$ form an edge in $H$.
Indeed, a graph $G$ is $k$-colourable if and only if there exists a homomorphism from $G$ into the complete graph on $k$ vertices, $K_k$.

Given a finite graph, there is an algorithm that decides whether it is 2-colourable or not in polynomial time in the size of the graph.
Informally, traversing the graph through its adjacencies while alternating colours produces either a valid 2-colouring, or a contradiction (the method reaches a vertex adjacent to an already-coloured vertex and is forced to use the same colour as its neighbour); by the limited number of colours, this contradiction implies that no 2-colouring exists.
For finite graphs, deciding $3$-colourability is NP-complete.

An elementary exercise for finite graphs proves that a graph is $2$-colourable if and only if it has no odd cycles.
There is a compactness argument that extends this result to infinite graphs, which is essentially the De Bruijn--Erd\H{o}s theorem~\cite[14.1.27]{bondy2008graph}.
This can be used to show that the \emph{infinite Hamming graph} $\mathcal H$, defined in the subsequent paragraph, is $2$-colourable.

Consider the Cantor space $2^\omega$ and impose a graph structure on it by declaring two infinite binary sequences $x$ and $y$ adjacent if they differ at exactly one coordinate.
The resulting graph, which we denote by $\mathcal H$, is $2$-colourable as it lacks any odd cycles.
Suppose that $A\subseteq 2^\omega$ is an independent set\footnote{An \emph{independent} set of vertices is a set disjoint from the edge relation, i.e.\ a set of pairwise non-adjacent vertices.}.
Suppose also that $A$ has the Baire property (for example, every Borel set has the Baire property).
If $A$ were not meagre, then it would be comeagre in some basic clopen set $[t]:=\{x\in2^\omega:t\subseteq x\}$.
Let $f\colon2^\omega\to2^\omega$ be the homeomorphism that flips the $(|t|+1)$-st\footnote{$|t|$ denotes the length of the finite sequence $t$, i.e.\ its domain.} coordinate, so that the image of $[t]$ under $f$ is $[t]$ itself, and every $x$ is adjacent to $f(x)$ in $\mathcal H$.
Then $f[A]$ is also comeagre in $[t]$, but $A\cap f[A]=\emptyset$ by independence, which contradicts the Baire category theorem applied to the Polish subspace $[t]$.
Therefore, every independent set with the Baire property is meagre, and thus $2^\omega$ cannot be covered by countably many such classes.

The preceding argument illustrates the cost of imposing descriptive (computable or measurable) requirements on the solution, in this case of the colouring problem: namely, that the number of colours must be increased.
These dichotomies share a common shape, formalized in descriptive set theory as the basis paradigm.

Kechris, Solecki, and Todor{\v{c}}evi{\'c}~\cite{kechris1999borel} introduced this notion of the \emph{Borel chromatic number} precisely to study this phenomenon.
Remarkably, they proved the following result, known as the $\mathbb{G}_0$ dichotomy: there exists a single Borel graph~$\mathbb{G}_0$ such that
\begin{align*}
    \chi_B(G) > \aleph_0\quad
    \Longleftrightarrow\quad
    \mathbb{G}_0 \to_B G.
\end{align*}

The symbol $\mathbb{G}_0 \to_B G$ abbreviates the phrase ``there is a Borel homomorphism from $\mathbb{G}_0$ to $G$.''
The graph~$\mathbb{G}_0$ forms a one-element \emph{basis} for the class of Borel graphs with uncountable Borel chromatic number.
One can interpret the existence of a small (i.e.\ singleton, or finite, or countable) basis as a sort of measure of complexity of the associated decision problem.
For instance, the existence of $\mathbb{G}_0$ says that deciding whether a Borel graph has uncountable Borel chromatic number is simple: one only has to verify one possible ill-behaved property, namely containing a homomorphic copy of $\mathbb{G}_0$, to give a positive answer; and in some sense, finding this homomorphic copy is the \emph{only} way of proving that a Borel graph has uncountable Borel chromatic number.
In contrast, deciding whether a Borel graph has infinite Borel chromatic number is a much more complex task, as no easily describable method exists for deciding this in terms of graph homomorphisms, and so different graphs probably require wildly different proof techniques, if any.

Similarly to the case for finite graphs in the opening paragraphs of the chapter, Carroy, Miller, Schrittesser, and Vidny{\'{a}}nszky~\cite{L0paper} proved the existence of a one-element basis for analytic graphs of Borel chromatic number at least three, that is, there exists a Borel graph $\mathbb{L}_0$ with the following property.
The result is stated as a corollary in Section~\ref{sec:concise_L0}.

\begin{restatable}[$\mathbb{L}_0$ dichotomy]{corollary}{Lzero}\label{cor:L0-dichotomy}
    For any analytic graph $G$ on a Polish space, $\chi_B(G)\geq 3$ if and only if there is a continuous mapping $\mathbb{L}_0\to_cG$.
\end{restatable}

\begin{remark}\label{rmk:finite-Borel}
    Finite graphs are considered trivial examples of analytic graphs.
    Concretely, for a finite space $X$ and a graph $G$ defined on it, $\chi_B(G)=\chi(G)$ and every set of vertices and every set of edges is both Borel and analytic.
    Note that if $G$ is finite, then $\chi(G)\geq3$ is equivalent to $G$ having an odd cycle;
    and it is not too difficult to show that this condition is equivalent to there being a continuous mapping $\mathbb{L}_0\to_cG$.
    In summary, the preceding result holds for finite graphs as well.
\end{remark}

A new proof of Corollary~\ref{cor:L0-dichotomy} is given in Section~\ref{sec:concise_L0}, and a new antibasis result for \emph{dicolourings}, a directed analogue of colourings, is established in Section~\ref{sec:dichromatic}.

It is worth seeing explicitly that the $\mathbb{G}_0$ dichotomy refines Theorem~\ref{thm:cantor-perfect}.
Let $C \subseteq 2^\omega$ be closed and consider the complete graph $K_C = (C \times C) \setminus \Delta_C$, which is closed and in particular analytic.
Apply the $\mathbb{G}_0$ dichotomy:
\begin{itemize}
    \item If $\chi_B(K_C) \leq \aleph_0$, then $C$ partitions into countably many $K_C$-independent Borel sets.
    But independent sets in a complete graph contain at most one point, so $C$ is countable.
    \item Otherwise, there is a continuous homomorphism $\phi\colon (2^\omega, \mathbb{G}_0) \to (C, K_C)$, that is, a continuous map $\phi\colon 2^\omega \to C$ sending $\mathbb{G}_0$-edges to pairs of distinct points.
    The closed equivalence relation $E_\phi = \{(x,y) : \phi(x) = \phi(y)\}$ is then disjoint from $\mathbb{G}_0$ off the diagonal, so each $E_\phi$-class is $\mathbb{G}_0$-independent, hence meagre; a standard Mycielski-style argument therefore yields a continuous injection $\iota\colon 2^\omega \to 2^\omega$ on which $E_\phi$ restricts to equality, and $\phi \circ \iota\colon 2^\omega \to C$ is the desired continuous injection.
\end{itemize}
Cantor's theorem thus appears as the instance of the $\mathbb{G}_0$ dichotomy for complete graphs on closed sets.

A recent paper by Todor{\v{c}}evi{\'c} and Vidny{\'{a}}nszky \cite{Todorcevic_2021_ComplexProbBorel} rules out the existence of countable bases for any other finite colour threshold.
In Section~\ref{sec:dichromatic}, a new result completes the picture for the \emph{dichromatic} number, an analogue of the chromatic number for directed graphs.

This is an illustrative introduction to the dichotomies that occupy the present thesis: the perfect set theorem which foreshadows the $\mathbb{G}_0$ dichotomy, and the Rosenthal-style dichotomies underlying the Deep Computations exposition of Chapter~\ref{ch:deep_computations}, as well as other complexity-theoretic themes that recur throughout this work.

The same dichotomy phenomenon, sequential vs.\ topological closure, Baire-1 vs.\ wild, drives Chapter~\ref{ch:deep_computations}, where we identify new definability classes for spaces of computations.

It is worth pausing to note that the basis paradigm where small obstructions characterize a structural property predates descriptive combinatorics by decades.
Two of the most celebrated results in graph theory take this form.
The following two theorems (see \cite[Section~10.5]{bondy2008graph}) are well-known characterizations of \emph{planarity} in graphs in these terms.
While the two bases consist of the same graphs, notice that the underlying ordering on graphs differs.

\begin{theorem}[Wagner, 1937]\label{thm:wagner-basis}
    The set $\{K_5, K_{3,3}\}$ is a basis for non-planarity under the minor ordering: a finite graph $G$ is non-planar if and only if there exists $B \in \{K_5, K_{3,3}\}$ such that $B \leq_{\mathrm{minor}} G$.
\end{theorem}

\begin{theorem}[Kuratowski, 1930]\label{thm:kuratowski-basis}
    The set $\{K_5, K_{3,3}\}$ is a basis for non-planarity under the topological minor ordering: a finite graph $G$ is non-planar if and only if there exists $B \in \{K_5, K_{3,3}\}$ such that $G$ contains a subdivision of~$B$.
\end{theorem}

Planar graphs are a particular example of a broader class of graphs for which a much more general finite basis result exists.

\begin{theorem}[Robertson--Seymour]\label{thm:robertson-seymour}
    Finite graphs are well-quasi-ordered under the minor relation.
    In particular, for every minor-closed class $\mathcal{C}$ of finite graphs, there exists a finite set $\mathcal{B}$ such that $G \notin \mathcal{C}$ if and only if $B \leq_{\mathrm{minor}} G$ for some $B \in \mathcal{B}$.
\end{theorem}

Nevertheless, there are classes of graphs with interesting results that follow a very similar theme for which the preceding theorem does not apply.
One such instance, discussed in Section~\ref{sec:minimal-rigid}, concerns asymmetric graphs, which are not subject to the Robertson--Seymour theorem by virtue of asymmetry failing to be minor-closed.

\begin{restatable}[Schweitzer--Schweitzer {\cite{SCHWEITZER2017215}}]{theorem}{schweitzerminimal}\label {thm:schweitzer-minimal}
	There exists a set $\mathcal{B}$ of exactly $18$ finite graphs, the minimal asymmetric graphs, such that every finite asymmetric graph $G$ contains some $B \in \mathcal{B}$ as an induced subgraph; that is, there is an injective homomorphism $\varphi\colon B \hookrightarrow G$ that preserves non-adjacency.
	Moreover, the elements of $\mathcal{B}$ are pairwise incomparable under the induced subgraph ordering.
\end{restatable}

A precursor to the study of the group action presented in Chapter~\ref{ch:group_actions} involves the following dichotomy, which is the main result of our article~\cite{Caroetal2023}.

\begin{theorem}[Part of Theorem~\ref{thm:constant_bal}]
    A graph $G$ with an even number of edges $k>0$ has constant balancing number\footnote{\emph{Constant balancing number} is defined in Subsection~\ref{subsec:amoebas_bal}.} if and only if there are injective homomorphisms $\frac{k}2K_2\to G$ and $K_{1,k/2}\to G$.
\end{theorem}